购买分区的日志不断增长，原结果、关闭的会话、未决付款、退款占额与事件记录也随之积累。第十七章把这些事实作为完整状态机复制，并在某个已应用索引处冻结图像。较大的图像需要适合大块字节的保存结构：传输中断后能够续传，单盘损坏后能够换副本读取，恢复者还能判断拿到的是哪一个完整版本。存储服务承载图像字节，购买分区继续决定业务事实。

这一子系统包含两条路径。元数据路径回答文件由哪些块组成、块在哪里、哪个版本已经发布。数据路径传输实际字节。若每次读一个大块都让中央元数据服务转发内容，它的网卡就会限制整个集群的带宽。客户端先取得位置，再直接访问数据节点，可以把大量传输分散到许多机器，同时保留明确的命名与发布入口。

下面沿三种保存接口看字节何时可读、失败后依据什么继续。GFS把并发修改和追加重试的结果交给应用按区间与记录解释，HDFS通过单写入者、流水线和块恢复收束未完成文件，对象存储则把完整对象与条件更新入口暴露给发布者。冻结图像适合先形成不可变内容，再发布固定集合；事件归档还可能采用持续追加。两类负载分别检验完整版本与重复记录，因而需要辨明各自的确认、读取和恢复边界。

\CNineteenFigure{ch19-01}{元数据决定命名与位置，数据节点直接传输图像字节}{fig:19-metadata}

\subsection{GFS怎样组织大文件}

Google File System，简称GFS，是为大规模数据处理设计的分布式文件系统。应用看见按路径命名的文件，文件系统把内容分散到许多普通机器上。它针对的负载以大文件、顺序扫描和持续追加为主，设计目标偏向许多客户端合起来的持续吞吐。机器和磁盘故障作为日常事件进入监测、补副本和恢复流程~\cite{gfs2003}。

因此，GFS没有把每个业务对象都做成一个独立小文件。归档器可以把许多事件写入一个大文件，分析作业连续读取这些事件；图像也可以作为较大的字节流保存。文件被切成称为chunk的大块，每个chunk有一个全局唯一且不变的句柄，句柄用于识别这个块。原设计的块大小为64 MB，默认保存三个副本。大块减少文件到块的映射条目，也让客户端在同一连接上连续读写更多数据。小文件或许只占一个chunk，大量客户端同时访问它时，少数副本又可能形成热点。块大小和副本数量因而与文件分布、访问局部性及重传成本共同决定容量安排。

集群中，master管理文件名、访问控制、文件到chunk的映射、租约和副本放置，chunkserver把chunk副本存为本地文件。客户端库把应用给出的文件名和字节偏移换算为块序号，向master查询该序号对应的句柄及副本地址，再选择一个副本直接读取指定字节范围。客户端缓存位置元数据，连续读取同一块时便可复用这些信息。

\CNineteenFigure{ch19-21-layout}{一个文件的有序块映射与块副本位置分别保存，三份副本共同保存同一个chunk}{fig:19-gfs-layout}

图~\ref{fig:19-gfs-layout} 中，文件的第0块与第1块对应不同句柄。同一个句柄的三个副本保存同一块的内容，不是文件中依次连接的三个区段。恢复文件先按文件映射确定块的次序，再为每一块选择可用副本。副本搬到另一台机器以后，句柄和文件内的块序号可以保持原样。

master把关键元数据变化写入操作日志，在本地和远端刷盘后答复相应元数据操作，并通过检查点与后续日志恢复命名空间。即使所有数据块都在，丢失文件到块的映射，也会失去可靠重建文件的依据。块的当前位置则主要从chunkserver启动报告和运行中的状态交流重建。磁盘坏了以后，中心保存的一条旧位置记录无法使块重新存在，由实际持有副本的节点报告状态更直接。

这种集中管理减少了全局放置和回收的协调复杂度，也把元数据规模、元数据操作速率及master恢复能力变成设计约束。大量字节绕过master，并没有消除命名空间的保存责任。副本还要跨共同故障域放置：三份都在同一机架时，电源或交换机故障会同时切断访问。补副本消耗正常读写使用的磁盘和网络，master限制节点及集群的并行复制，避免恢复流量耗尽前台容量。

\subsection{写入的数据路径与修改顺序}

对一个chunk的修改需要各副本采用相同次序。master给一个副本授予有期限的primary lease，即主副本租约。在租约有效期间，这个primary为修改编号，其他副本采用它给出的顺序。租约授予次序与租约内的修改序号共同确定修改次序。客户端缓存primary的位置，访问失败或租约失效后再询问master。

大块数据先推送到各副本的缓冲区，推送链可以按网络距离选择。等所有副本确认收到数据，客户端才向primary发出修改请求，引用先前推送的数据。primary选择修改序号并本地应用，再通知其他副本按同一序号应用。其他副本完成后回复primary，primary再回复客户端。数据经过哪条传输链，与谁决定修改次序，分别适合按带宽和一致顺序优化。

\CNineteenFigure{ch19-02}{先把数据送达副本，再由primary确定同一修改次序}{fig:19-gfs-write}

如果某个副本在修改时失败，primary和其他一些副本可能已经写入。客户端得到失败后重试，系统面对的是局部完成的历史。GFS为这类大文件负载规定了较宽松的一致性模型，相关文献明确区分成功串行写、并发成功写和失败写影响的文件区间~\cite{gfs2003}\footnote{这里的一致区间表示从不同副本读取会得到相同数据，已定义区间还要求完整体现相应修改的内容。完整定义与应用处理办法见集中书目中GFS的核读范围。}。

一次没有并发干扰的成功写，使所影响区间包含完整写入内容。应用的一次大写入若被拆成多次底层修改，另一客户端可以在它们之间写入，所有副本即使采用同一顺序，最终区间也可能混合两方内容。修改失败时，部分副本已经改变、另一些尚未改变，读取还可能因副本不同而不同。文件名创建的原子性、单块修改的顺序与应用整份内容的完整发布，需要分别从对应记录判断。

\CNineteenFigure{ch19-22-regions}{串行成功写、交错的大写入和部分失败，对文件区间形成不同结果}{fig:19-gfs-regions}

master在授予新租约前推进chunk版本，以识别错过修改的旧副本。旧副本不参与当前写入，也不作为当前位置查询的结果返回。客户端已经缓存的旧地址可能暂时指向旧副本，缓存到期、重新打开文件或读到过早的文件尾以后，再查询有效位置。恢复读取因此还受到位置缓存的有效期约束。

\subsection{追加记录与读取者的责任}

GFS的record append由客户端提供记录，系统选择追加位置。若当前chunk放不下记录，primary可以填充剩余空间，让客户端到下一个chunk重试。一次成功追加保证记录至少一次作为连续字节区间保存，并返回该区间的起始偏移。失败后的重试可能留下重复记录、片段或填充，副本上的失败区间也可能不同~\cite{gfs2003}。

例如，归档器把\code{EV1}及完整payload追加到文件。第一份记录已经写成，成功回复丢失，归档器又追加同一记录。第二次的位置可能不同，文件中也可能有两份完整\code{EV1}。读取者先用长度、格式版本与内容校验排除片段和填充，再按事件身份与原payload核对重复。同一\code{EV1}携带不同payload时，进入冲突处理。

\CNineteenFigure{ch19-03}{追加重试可能形成第二份完整记录，消费方按原事件关联已受理效果}{fig:19-append}

第十六章的履约消费者在一个事务中保存Inbox记录I、业务守卫、实际履约编号F与dispatch意图。相同事件重投时，它恢复这份原关联。归档文件保存了几份记录，描述的是字节保存历史；是否已接受履约，由消费保存域的I、业务守卫与F共同回答。

GFS还为chunk内的64 KB区段保存校验和。chunkserver读出前检验本地数据，校验不匹配时返回错误并报告master，读取者换副本，master从有效副本补出新副本。record append允许部分失败区间不同，因此不能用两份chunk的全部字节相等代替逐副本的损坏检测。校验确认的是字节与校验信息匹配，图像来自哪个业务前缀、是否获准恢复，还要由来源与发布清单说明。

\subsection{HDFS的命名空间、数据流水线与恢复}

Hadoop Distributed File System，简称HDFS，同样服务于大数据集和流式访问。应用通过文件路径访问数据，文件被切成块，由多个DataNode保存副本。它以写入后多次读取为常见用法，支持追加和截断，一个打开写入的文件在同一时刻只有一个写入者。限定写入者和修改方式，可以把大量并发读取与数据传输留给数据节点，并简化同一文件写入的协调~\cite{hdfs342-storage}。

NameNode管理目录、文件、文件到块的关系与副本位置，DataNode保存块并直接服务客户端。NameNode把命名空间变化写入EditLog，FsImage保存一个已经形成的命名空间图像，启动时结合图像与后续日志恢复。DataNode通过心跳报告存活情况，通过块报告提供持有的块。命名空间图像是文件系统的元数据图像，本书购买分区的完整图像则是由它承载的一份文件内容。

\CNineteenFigure{ch19-23-hdfs}{NameNode保存命名空间，DataNode流水线保存块，确认沿相反方向返回}{fig:19-hdfs}

客户端创建或追加文件时，先取得写入权限和目标块的副本位置，再把数据包送入DataNode流水线。第一个节点接收并向后转发，后续节点依次保存数据，确认沿链返回。流水线中断后，客户端根据已确认位置、存活副本和恢复协议重建流水线。写入方法返回、其他读取者看见数据、持久同步完成和文件关闭，各自对应接口及配置规定的进度。

写入者尚未关闭文件就退出，写入租约最终触发恢复。NameNode推进恢复代次，选定的DataNode收集末块各副本的generation stamp、长度与状态。generation stamp标记块的写入或恢复代次，用于排除旧写入。恢复者据副本状态选择可接受的共同结果，同步参与副本，再向NameNode提交块与文件的更新。

长度选择依赖状态。Hadoop 3.4.2的连续块恢复分支发现FINALIZED副本时，核对这些已完成副本的长度一致，并按该长度选择参与副本。若最佳状态是正在写入的RBW或等待恢复的RWR，则在最佳状态集合中取最小长度。只到达某个副本的更长尾部，可能尚未成为可共同保存的前缀~\cite{hdfs342-storage}。

\CNineteenFigure{ch19-04}{租约恢复先核对代次、状态与长度，再结束未完成写入}{fig:19-hdfs-recovery}

高可用NameNode把元数据日志放在JournalNode组，由多数JournalNode接受EditLog写入。Standby跟随日志，接替前追上所需状态。旧Active失去写日志资格后，短时间内还可能提供过时读取，部署中的隔离和客户端故障转移共同处理这条路径。Secondary NameNode主要承担检查点工作，名称本身不意味着它自动接替Active。

DataNode保存并传输数据的校验信息，读取者发现损坏时换副本。块编号、generation stamp和租约对应文件及块的保存责任。用于购买图像时，发布者还要比较恢复后文件的长度和整幅摘要：租约恢复可能舍去未共同确认的末尾，恢复后的合法文件因此可能短于业务期望的图像。导入器依据下一节的manifest，判断它是否已经取得完整业务前缀。

\subsection{对象存储的对象、版本与发布入口}

对象存储把一串字节及其属性放到对象键下。调用方通常在一个bucket中用键定位对象，通过API创建、读取、删除或条件更新。bucket可以理解为提供命名和访问管理边界的对象容器。文件系统把路径、目录和块向客户端组合为文件，对象接口则以完整对象为基本访问单位，内部怎样复制字节由服务实现。

这适合把图像切成不可变对象，再用一份小清单连接它们。应用可以并行上传不同块，也可以重传一个失败块。代价是跨多个对象的完整发布由应用明确组织，一个对象创建成功，只能证明这个对象取得了相应保存结果。对象版本、访问权限、生命周期与读取一致性，则按所用服务的公开接口配置~\cite{s3-object-api}。

以S3分段上传为例，客户端先创建上传，取得\code{upload\_id}，再按part number上传各部分。完成请求列出要组合的部分及ETag，服务按部分编号顺序形成对象。ETag是接口返回的标识，可用于相应条件比较，分段上传或加密等情形下不能一概把它视作整个文件的MD5。应用另外保存完整图像的期望长度和摘要，核对完成后的内容。

Complete请求还可能在初始HTTP 200之后，于响应体中报告错误。直接调用API的客户端因此解析完整响应，依据最终结果继续。重试完成请求、查询对象和重传数据，各自沿原上传身份与对象键恢复，避免仅凭HTTP状态行提前发布~\cite{s3-object-api}。

创建对象时，\code{If-None-Match: *}表达键当前不存在的前提。更新发布指针时，\code{If-Match}表达所读ETag依然匹配的前提。条件失败后，发布者重读入口，判断自己的图像是否已发布，或是否有另一发布者先完成。版本化bucket中的version ID与ETag各有含义，清单按所采用的接口固定确切对象或对象版本。

\CNineteenFigure{ch19-05}{分段上传形成完整对象，条件更新另外决定图像的发布入口}{fig:19-object}

恢复者从一份确定清单进入，只读取清单列出的块。这使大数据传输和小元数据发布有了清楚的先后：先把全部必要字节保存到要求的副本或对象，再让入口引用这个完整集合。
\FloatBarrier
